% Chapter 2
\chapter{Literature Review}
\label{Chapter2}
\lhead{Chapter 2. \emph{Literature Review}}

\section{Introduction}
This chapter reviews relevant work on magnetic field monitoring, anomaly detection in time series, recurrent neural networks for prediction, and direction finding from tri-axial magnetometer data. It provides the background for the design choices adopted in Magnavis.

\section{Magnetic Anomaly Detection}
Magnetic anomaly detection aims to identify short-lived deviations from the ambient or expected magnetic field. In geophysical and defence applications, such anomalies can indicate buried ferrous objects, UXO, or other sources. Traditional methods include threshold-based detection, statistical outlier detection (e.g. based on mean and standard deviation or interquartile range), and model-based approaches where a baseline is estimated and residuals are analysed. Magnavis follows a model-based approach: a GRU predictor learns the baseline, and anomalies are flagged when the prediction error exceeds an adaptive threshold.

\section{Recurrent Models for Time-Series Prediction}
Recurrent neural networks (RNNs), including Long Short-Term Memory (LSTM) and Gated Recurrent Unit (GRU) architectures, are widely used for time-series forecasting. They can capture temporal dependencies and are suitable for learning smooth baselines such as diurnal magnetic variation. The Magnavis predictor uses a GRU layer (32 units) followed by dense layers, with cyclic time-of-day (and optionally yearly) features to improve generalisation. Pre-training on extended historical data (e.g. 2--4 months) is supported for better accuracy and faster startup.

\section{Observatory-Based Magnetic Monitoring}
Multi-station magnetic observatories provide spatially distributed measurements for studying geomagnetic variations and local anomalies. When observatories are equipped with tri-axial magnetometers, the three components (e.g. North, East, Vertical or sensor-frame components) can be used to compute the direction of the magnetic field perturbation. The direction of the perturbation vector is often used as a first approximation to the bearing of the source, though it is not in general equal to the direction to the source; dipole fitting or triangulation can improve localisation.

\section{Direction Finding and Triangulation}
Direction finding from tri-axial data typically involves: (i) defining a baseline (e.g. median over a quiet period) for each component, (ii) computing the perturbation vector as the difference from the baseline, (iii) deriving azimuth and inclination from the horizontal and vertical components. Triangulation uses bearing lines from two or more observatories; their intersection (or closest approach in 3D) gives an estimated source position. Accuracy depends on baseline length, crossing angle, and time alignment of anomalies across sites. Alternative methods include dipole fitting (inverse problem) and multi-point dipole inversion, which use the full vector and magnitude at one or more points.

\section{Anomaly Detection in Time Series}
Statistical anomaly detection in time series often uses moving statistics (mean, median, standard deviation) and flags points that exceed a multiple of the standard deviation from the mean. Magnavis uses an adaptive threshold (mean of prediction errors plus a multiplier times the standard deviation) and an optional freeze window: after the first anomaly, the threshold can be frozen for a configurable duration to avoid rapid re-triggering. Interpolation of predicted values at actual timestamps ensures accurate comparison and reduces timing errors.

\section{Summary}
The reviewed literature supports the choice of a GRU-based predictor for baseline learning, an adaptive threshold on prediction residuals for anomaly detection, and perturbation-vector-based direction finding with optional triangulation. The Magnavis framework integrates these elements in a single application and extends them with multi-sensor support, pre-trained models, and a DB-backed workflow.
